\section{Schema-first bounded reflection}
\label{sec:reflection}

\subsection{Why a schema rather than a term compiler}

A term-first bounded-formula compiler takes a syntax tree \(A\) with free
variables and returns a term \(\chi_A\).  If the variables are later replaced
by compound terms, one must compare two constructions:
\[
   \chi_{A[\sigma]}
   \qquad\text{and}\qquad
   \chi_A[\sigma].
\]
At bounded binders these may contain different primitive-recursive symbols,
because closure conversion has been performed before or after substitution.
They compute the same function, but proving that equality requires another
induction through the compiler and through the finite recursion.

\Scope removes this obligation by changing the interface.  A bounded formula
is a \emph{scheme} over \(k\) formal parameters.  Compilation returns one fixed
symbol
\[
   \chi_A : \PR\,k.
\]
For arbitrary terms \(\vec t\), the characteristic term is always
\(\chi_A(\vec t)\).  There is no operation that recompiles the scheme after
substitution.  Instantiation of the formula and instantiation of the bit both
use the same vector \(\vec t\).

\subsection{Reflected atomic schemes}

\begin{definition}[Atomic schema]
A reflected atomic scheme of arity \(k\) consists of
\begin{enumerate}[label=(\alph*)]
\item a formula \(A\in\Fm\,\PTm\,k\);
\item a symbol \(b_A\in\PR\,k\);
\item a closed derivation of Booleanity
\[
  \turnstilep b_A(\vec x)=\zero\ \vee\ b_A(\vec x)=\one;
\]
\item a closed derivation
\[
  \turnstilep\bigl(b_A(\vec x)=\one\bigr)\leftrightarrow A(\vec x).
\]
\end{enumerate}
\end{definition}

Equality and order atoms are obtained from the checked arithmetic library.  For
example, \code{eqb(s,t)} is \(1\) when \(s=t\) and \(0\) otherwise; the strict
order formula is an equality involving truncated subtraction and can therefore
be reflected by the same equality bit.  A theorem packet may also introduce a
new atom by supplying its own fixed PR verifier and object derivations.

\begin{lemma}[Falsity reflection from Booleanity]
\label{lem:atomic-zero}
For every reflected atom,
\[
  \turnstilep
  \bigl(b_A(\vec x)=\zero\bigr)\leftrightarrow\neg A(\vec x).
\]
\end{lemma}

\begin{proof}
If \(b_A=0\) and \(A\), truth reflection gives \(b_A=1\), contradicting
\(0\neq1\).  Conversely assume \(\neg A\).  Booleanity gives
\(b_A=0\vee b_A=1\).  The first disjunct is the goal; the second yields \(A\)
by truth reflection and hence absurdity.  This is an intuitionistic proof by an
object-level finite disjunction, not an appeal to unrestricted excluded middle.
\end{proof}

\subsection{Bounded formula schemes}

For each \(k\), let \(\mathsf{BCode}(k)\) be generated by
\[
\begin{array}{c}
  \mathsf{atom}(a),\quad A\wedge B,\quad A\vee B,\quad A\to B,
  \quad\neg A,\\[0.3em]
  \exists i<t\;A(i,\vec x),
  \qquad
  \forall i<t\;A(i,\vec x).
\end{array}
\]
In a bounded body the newest variable \(0\) is \(i\), while the old parameters
are shifted.  The human-facing denotation is the ordinary first-order formula
\[
\begin{aligned}
 \llbracket\exists i<t\ A\rrbracket
   &= \exists i\,(i<t\wedge\llbracket A\rrbracket),\\
 \llbracket\forall i<t\ A\rrbracket
   &= \forall i\,(i<t\to\llbracket A\rrbracket).
\end{aligned}
\]
No claim is made that these formulas are quantifier-free.  Only their
characteristic-bit equations are used as induction invariants.

\subsection{Arity-generic bounded search}

The checked repository search accepts one bound and one packed parameter.  The
underlying proof is already arity-independent: \code{recDef} permits any
number of parameters.  We use the following direct generalization.

\begin{theorem}[Parameterized least-zero search]
\label{thm:searchN}
Let \(p(i,\vec x)\) be a fixed primitive-recursive term.  There is a fixed
symbol \(\mu_p(N,\vec x)\) such that iPRA proves
\begin{align*}
 &\mu_p(N,\vec x)\le N,\tag{S1}\\
 &\mu_p(N,\vec x)<N
      \longrightarrow p(\mu_p(N,\vec x),\vec x)=0,\tag{S2}\\
 &i<\mu_p(N,\vec x)
      \longrightarrow p(i,\vec x)\neq0.\tag{S3}
\end{align*}
Moreover, if \(i<N\) and \(p(i,\vec x)=0\), then
\[
  \mu_p(N,\vec x)<N,\qquad
  p(\mu_p(N,\vec x),\vec x)=0,
  \qquad
  \mu_p(N,\vec x)\le i.
\]
\end{theorem}

\begin{proof}
Use the same primitive recursion as \code{Core/Search.lean}: at stage \(N+1\),
retain the old candidate unless it equals \(N\) and \(p(N,\vec x)=0\).  The
joint invariant is exactly (S1)--(S3), written with the current stage, a test
index, and the parameters as free variables.  It is a conjunction of equality,
order, implication, and negated equality, hence syntactically quantifier-free.
Quantifier-free induction on \(N\) proves the invariant.  The found theorem is
then a finite order argument: if the least candidate were after a supplied
zero, (S3) would contradict that zero.  This is precisely the existing proof
with the tuple \(\vec x\) left as separate parameters.  In Lean,
\code{PTm.toPR} closes the step term and \code{recDef} supplies the arbitrary
arity.
\end{proof}

\subsection{Compilation}

Assume fixed primitive-recursive bit operations
\(\mathsf{notb},\mathsf{andb},\mathsf{orb},\mathsf{impb}\), each with object
proofs on Boolean inputs.  Define \(\chi_A\) recursively.

\begin{align*}
 \chi_{\neg A} &= \mathsf{notb}(\chi_A),\\
 \chi_{A\wedge B} &= \mathsf{andb}(\chi_A,\chi_B),\\
 \chi_{A\vee B} &= \mathsf{orb}(\chi_A,\chi_B),\\
 \chi_{A\to B} &= \mathsf{orb}(\mathsf{notb}(\chi_A),\chi_B).
\end{align*}
For a bounded existential, put
\[
  p_A(i,\vec x)=\mathsf{notb}(\chi_A(i,\vec x)),
  \qquad m=\mu_{p_A}(t(\vec x),\vec x),
\]
and let \(\chi_{\exists i<t A}\) be the characteristic bit of \(m<t\).
For a bounded universal, search directly for a zero of \(\chi_A\):
\[
  r=\mu_{\chi_A}(t(\vec x),\vec x),
  \qquad
  \chi_{\forall i<t A}=\mathsf{eqb}(r,t(\vec x)).
\]
All of these are fixed symbols of arity \(k\), built by composition from the
body symbol, the bound's \code{PTm.toPR} symbol, and parameterized search.

\begin{theorem}[Uniform bounded reflection]
\label{thm:reflection}
For every fixed \(A\in\mathsf{BCode}(k)\), the compiler effectively constructs
\(\chi_A\in\PR\,k\) and closed iPRA derivations
\begin{align*}
 &\turnstilep \chi_A(\vec x)=0\ \vee\ \chi_A(\vec x)=1,
       \tag{R0}\\
 &\turnstilep
   \bigl(\chi_A(\vec x)=1\bigr)
       \leftrightarrow \llbracket A\rrbracket(\vec x),
       \tag{R1}\\
 &\turnstilep
   \bigl(\chi_A(\vec x)=0\bigr)
       \leftrightarrow \neg\llbracket A\rrbracket(\vec x).
       \tag{R2}
\end{align*}
The construction is primitive recursive at the metalevel and introduces no
object-language syntax evaluator.
\end{theorem}

The proof occupies the next two sections.  The important structural fact is
that the output of compilation is a \(\PR\) symbol.  Once (R0)--(R2) are
proved at formal variables, the checked \code{inst} theorem specializes them
under arbitrary contexts and arbitrary object terms.
